\section{Guaranteed performance values}

\subsection{Typical vs.\ guaranteed}
Indicative ranges from the CIGRE biennial HVDC reliability surveys, together with
values that an OEM can realistically stand behind:

\begin{center}
\renewcommand{\arraystretch}{1.35}
\begin{tabularx}{\linewidth}{@{}l X c c@{}}
  \toprule
  \textbf{Metric} & \textbf{Description} & \textbf{Typical} & \textbf{Guaranteed} \\
  \midrule
  FEU & Forced energy unavailability    & 0.2--1.0\,\%      & $\approx$0.5\,\% \\
  SEU & Scheduled energy unavailability & 1--3\,\%          & $\approx$1.0\,\% \\
  FOR & Forced outage rate (pole)       & 4--8 / yr         & 2--3 / yr \\
      & Forced outage rate (bipole)     & $\lesssim$0.5 / yr & $\approx$0.1 / yr \\
  \bottomrule
\end{tabularx}
\end{center}

\subsection{Interpretation}
\begin{itemize}\setlength{\itemsep}{0.25em}
  \item Actual figures vary with technology (LCC vs.\ VSC), scheme age and
        transmission medium: long overhead lines raise forced-outage counts,
        cable schemes lower them.
  \item SEU is driven mainly by valve-maintenance duration and can be reduced by
        multi-shift working during outages.
  \item Very aggressive bipole figures (e.g.\ 0.1 outages/year, or 1 in 20 years)
        have been \emph{guaranteed on some projects but are commercially driven,
        not technically justified}; the residual exposure is effectively carried
        as commercial risk through the LD regime.
  \item VSC bipolar schemes with a dedicated metallic return are still rare, so
        CIGRE survey data for them is limited --- guarantees carry more
        uncertainty.
\end{itemize}

\subsection{RTEi availability-planning inputs}
Additional availability-planning data provided by RTE International (RTEi) gives
configuration-specific forced-outage rates and repair times:

\begin{center}
\renewcommand{\arraystretch}{1.35}
\begin{tabularx}{\linewidth}{@{}l X c c@{}}
  \toprule
  \textbf{Configuration} & \textbf{Outage type} & \textbf{FOR} & \textbf{MTTR} \\
  \midrule
  SMP            & Station / pole outage & 0.8 / (yr$\cdot$station) & 4--24\,h \\
  BIP with DMR   & Pole outage           & 0.8 / (yr$\cdot$station) & 4--24\,h \\
  BIP with DMR   & Bipole (full) outage  & 0.04--0.06 / yr          & --- \\
  \bottomrule
\end{tabularx}
\end{center}

\begin{itemize}\setlength{\itemsep}{0.25em}
  \item \textbf{Pole/station outages are frequent but shallow.} At
        $\approx$0.8 per year and station, a pole or single-station forced outage
        is expected roughly once a year, but with a short repair time
        (MTTR 4--24\,h) and, for a bipole, only a \emph{partial} (single-pole)
        loss of capacity.
  \item \textbf{Bipole (full) outages are rare.} A forced bipole outage rate of
        0.04--0.06 per year corresponds to \textbf{one event every 15--25
        years} --- more optimistic than the $\approx$0.1/yr ($\approx$1 in 10
        years) figure used elsewhere in this report, and consistent with the DMR
        configuration's redundancy.
  \item \textbf{Implication for LDs.} The high pole-outage frequency drives FOR
        exposure and repair logistics, while the bipole full-outage rate governs
        the deep-unavailability events that dominate the data-center availability
        case. Both should be reflected separately in the guarantee and LD
        structure.
\end{itemize}

{\color{atGray}\footnotesize Source: RTE International (RTEi) availability
planning.}

\subsection{Monitoring period}
\begin{itemize}\setlength{\itemsep}{0.2em}
  \item A monitoring period of \textbf{3 years} is typical.
  \item A robust buyer approach (common in North American projects) is to
        \textbf{extend to 4 years and count the best 3 of 4}, protecting against
        a single anomalous year.
  \item A simple mean/median-over-the-period approach is generally too risky for
        the buyer and should be avoided.
\end{itemize}
